% ID: SEL-D0065-OA-E04
% ESTADO: cotejada. Segunda pasada visual y borrador compilado revisados.
% Fuente: Selectividad/sel_2010_matematicas_aplicadas/mat_apli/mat_apli_sept/mat_apli_sept_ra/EXAMEN-ANDALUCÍA-1.pdf
% SHA-256: 35f05b19d43a5c26840240ff45ede93938ac575e4378437433e833cac887a341
% Página: 1; opción A; ejercicio original 4.

\begin{enumerate}[label=\alph*)]
\item \textbf{(1.25 puntos)} En una población de 2000 hombres y 2500 mujeres se quiere seleccionar una muestra de 135 personas mediante muestreo aleatorio estratificado con afijación proporcional, \textquestiondown{}cuál sería la composición de la muestra?
\item \textbf{(1.25 puntos)} Dada la población $\{6,8,11,a\}$, \textquestiondown{}cuánto debe valer $a$ sabiendo que la media de las medias muestrales de tamaño 3, obtenidas mediante muestreo aleatorio simple, es $10.3$?
\end{enumerate}
